\documentclass[11pt]{article}
\usepackage{amsmath,amssymb,amsthm}
\usepackage[margin=1.1in]{geometry}
\usepackage{graphicx}

\newtheorem{theorem}{Theorem}
\newtheorem{definition}{Definition}
\newtheorem{remark}{Remark}

\title{Minimization and Maximization Problems:\\
When Does Setting the Gradient to Zero Work?}
\author{Yuval Lavie}
\date{\today}

\begin{document}
\maketitle

\noindent
Every MLE note in this repository solves $\nabla(\cdot) = 0$ at some
point. This page states precisely when that move is justified, what can go
wrong, and what replaces it under constraints (Lagrange, KKT).
\texttt{optimization\_demo.py} demonstrates the failure modes and checks
the constrained solutions numerically.

\section{The objects: derivative, gradient, subgradient, Hessian}

Four objects carry all the local information optimization ever uses ---
the first two describe slope, the third rescues slope at kinks, the
fourth describes curvature.

\begin{definition}[Derivative --- three equivalent forms]
For $f:\mathbb{R}\to\mathbb{R}$, when the limit exists,
\begin{equation}
  f'(x)
  \;=\; \underbrace{\lim_{h\to 0}\frac{f(x+h)-f(x)}{h}}_{\text{(i)
  difference quotient}}
  \;=\; \underbrace{\lim_{y\to x}\frac{f(y)-f(x)}{y-x}}_{\text{(ii)
  two-point form}},
  \label{eq:derivative}
\end{equation}
and (iii) the \emph{linear-approximation form}: $f'(x)$ is the unique
number $a$ with
\begin{equation}
  f(x+h) \;=\; f(x) + a\,h + o(h)
  \qquad (\text{error} \to 0 \text{ faster than } h).
  \label{eq:derivlittleo}
\end{equation}
(i) and (ii) are the same statement under $y = x+h$; (iii) is the one
that \emph{generalizes} --- replace $a\,h$ by $\nabla f(x)^\top h$ and
it defines the gradient, by $J\,h$ and it defines the Jacobian. Rules
for computing are the derivations note (\texttt{rules.tex}); this page
needs what it means: the instantaneous trade rate between moving and
gaining. Figure~\ref{fig:deriv} (left) draws definition (i): secant
slopes collapsing onto the tangent as $h \to 0$.
\end{definition}

\begin{remark}[One-sided and symmetric quotients]
Two relatives of \eqref{eq:derivative} matter. The \emph{one-sided}
derivatives $f'(x^+)$ and $f'(x^-)$ take the limit through $h>0$ or
$h<0$ only; the derivative exists iff both exist and agree, and at a
kink they disagree ($|x|$ at $0$: $+1$ vs $-1$) --- exactly the gap the
subgradient below will fill. The \emph{symmetric} (central) quotient,
\begin{equation}
  \frac{f(x+h)-f(x-h)}{2h}
  \;=\; \tfrac12\Bigl[\text{forward} + \text{backward}\Bigr],
  \label{eq:symmetric}
\end{equation}
is \emph{not} an equivalent definition --- it is better and worse at
once. Better: when $f'(x)$ exists, \eqref{eq:symmetric} converges to it
with error $O(h^2)$ instead of the forward quotient's $O(h)$ (odd-order
terms cancel), which is why numerical gradient checks use central
differences. Worse: it can converge where no derivative exists --- on
$|x|$ at $0$ it returns $0$ for \emph{every} $h$ (fig.~\ref{fig:deriv},
right), averaging the $\pm1$ slopes into an answer that describes
nothing about the kink. Existence of the symmetric limit does not
imply differentiability.
\end{remark}

\begin{figure}[h]
  \centering
  \includegraphics[width=\linewidth]{fig_derivative.pdf}
  \caption{Left: the derivative as the limit of secant slopes
    (definition (i)). Right: the symmetric quotient's failure mode ---
    on $|x|$ at $0$ it converges to $0$ although no derivative exists
    (\texttt{optimization\_figures.py}).}
  \label{fig:deriv}
\end{figure}

\begin{definition}[Gradient]
For $f:\mathbb{R}^n\to\mathbb{R}$, the gradient collects the partial
derivatives,
\begin{equation}
  \nabla f(x) =
  \Bigl(\tfrac{\partial f}{\partial x_1},\dots,
        \tfrac{\partial f}{\partial x_n}\Bigr)^{\!\top},
  \qquad
  f(x+h) \approx f(x) + \nabla f(x)^\top h .
  \label{eq:gradient}
\end{equation}
Two geometric facts used constantly: $\nabla f(x)$ points in the
direction of \emph{steepest ascent} (so $-\nabla f$ is gradient
descent's direction), and $\nabla f$ is \emph{perpendicular to the
level set of $f$ through $x$} --- the set
$\{y : f(y) = f(x)\}$ of all points sharing $x$'s value (a contour
line): for any direction $d$ tangent to it, $\nabla f(x)^\top d = 0$,
since moving along the contour changes $f$ by nothing to first order.
This perpendicularity is the fact the Lagrange section runs on.
\end{definition}

\begin{definition}[Subgradient]
Where convex $f$ has a kink, no single slope exists, but many lines
touch from below: $g$ is a \emph{subgradient} at $x$ if
\begin{equation}
  f(y) \;\ge\; f(x) + g^\top (y - x) \qquad \text{for all } y,
  \label{eq:subgradient}
\end{equation}
and the set of all such $g$ is the subdifferential $\partial f(x)$.
At differentiable points $\partial f(x) = \{\nabla f(x)\}$ --- exactly
one line touches from below, the tangent (fig.~\ref{fig:subgrad},
left). At a kink a whole fan touches: for convex $f$ on $\mathbb{R}$,
\begin{equation}
  \partial f(x) = \bigl[\,f'(x^-),\; f'(x^+)\,\bigr],
  \label{eq:subdiffinterval}
\end{equation}
every slope between the one-sided derivatives --- for $|x|$ at $0$,
$\partial f(0) = [-1,1]$ (fig.~\ref{fig:subgrad}, right; the extreme
members $\pm1$ are the arms of $|x|$ themselves). Optimality becomes
$0 \in \partial f(x)$ --- the flat line is among the touching lines.
This is the object that makes ReLU, hinge, and $\ell_1$ losses
optimizable at all.
\end{definition}

\begin{figure}[h]
  \centering
  \includegraphics[width=\linewidth]{fig_subgradient.pdf}
  \caption{Subgradients as lines touching from below: a single tangent
    at a smooth point, the fan $[-1,1]$ at the kink of $|x|$
    (\texttt{optimization\_figures.py}).}
  \label{fig:subgrad}
\end{figure}

\begin{definition}[Hessian]
The matrix of second partials,
\begin{equation}
  \bigl(\nabla^2 f(x)\bigr)_{ij}
  = \frac{\partial^2 f}{\partial x_i\,\partial x_j},
  \qquad
  f(x+h) \approx f(x) + \nabla f(x)^\top h
  + \tfrac12\, h^\top \nabla^2 f(x)\, h ,
  \label{eq:hessian}
\end{equation}
symmetric under mild regularity (Schwarz). The gradient says which way
is up; the Hessian says how the slope itself changes --- curvature. Its
definiteness is the local shape: $\nabla^2 f \succ 0$ bowl,
$\prec 0$ dome, indefinite saddle; its eigenvalues are the curvatures
along its eigenvectors, and their spread is the conditioning that
Section 3's $\kappa$ quantifies.
\end{definition}

\section{Unconstrained problems: stationarity}

Maximization is minimization of the negation, so we speak of
$\min_\theta f(\theta)$ throughout.

\begin{definition}[Stationary point]
$\theta^\ast$ is \emph{stationary} for differentiable $f$ if
$\nabla f(\theta^\ast) = 0$.
\end{definition}

The meaning, in one breath: at a stationary point the linear
approximation \eqref{eq:gradient} is \emph{flat} --- to first order,
moving in any direction changes nothing, the function is momentarily
standing still (hence the name). No direction of first-order
improvement exists, which is why every interior optimum must be
stationary; but ``standing still'' is symmetric --- valley floors,
hilltops, and mountain passes all qualify --- so stationarity alone
cannot tell a minimum from a maximum from a saddle. That verdict
belongs to the curvature: the Hessian's definiteness
\eqref{eq:hessian} classifies the stationary point.

\begin{theorem}[First-order necessary condition, Fermat]
If $\theta^\ast$ is a local minimizer of differentiable $f$ and lies in
the \emph{interior} of the domain, then
\begin{equation}
  \nabla f(\theta^\ast) = 0 .
  \label{eq:stationarity}
\end{equation}
\end{theorem}

Necessary --- not sufficient. The catalogue of failures:
\begin{itemize}
  \item \textbf{Saddles and maxima.} $f(x,y)=x^2-y^2$ has
    $\nabla f = 0$ at the origin, which is neither (script, panel 1).
    Second-order check: $\nabla^2 f \succ 0$ at a stationary point
    $\Rightarrow$ local min; indefinite $\Rightarrow$ saddle.
  \item \textbf{Boundary optima.} On a constrained domain the minimum can
    sit on the boundary with $\nabla f \ne 0$ (Sections 4--5 handle this
    properly).
  \item \textbf{Non-differentiability.} $f(x)=|x|$ minimizes at $0$ with
    no derivative there; subgradients ($0 \in \partial f$) generalize
    \eqref{eq:stationarity}.
  \item \textbf{Solvable $\ne$ solvable in closed form.} For logistic
    regression, $\nabla = 0$ has no algebraic solution --- the equation is
    transcendental. That is not a failure of stationarity but of algebra;
    it is \emph{why iterative methods exist} (cf.\ the optimizers page).
    Our Bernoulli/Gaussian MLEs are the lucky cases where the stationary
    equations untangle.
\end{itemize}

\section{Convexity: when stationarity is the whole story}

\begin{definition}
$f$ is \emph{convex} if
$f(\lambda u + (1-\lambda)v) \le \lambda f(u) + (1-\lambda) f(v)$ for all
$u,v$ and $\lambda\in[0,1]$. Differentiable characterizations: $f$ convex
$\iff f(v) \ge f(u) + \nabla f(u)^\top (v-u)$ (tangents underestimate)
$\iff \nabla^2 f \succeq 0$.
\end{definition}

\begin{theorem}
For convex differentiable $f$, every stationary point is a \emph{global}
minimizer; if $f$ is strictly convex, the minimizer is unique.
\end{theorem}

This theorem is the license behind every ``set the derivative to zero''
in this repository: the Bernoulli log-likelihood (strictly concave), the
Gaussian NLL in each coordinate, least squares with full-rank $X$.
Two quantitative grades refine convexity, and between them govern
first-order convergence rates:
\begin{equation}
  \underbrace{f(v) \ge f(u) + \nabla f(u)^\top(v-u)
  + \tfrac{\mu}{2}\lVert v-u\rVert^2}_{\text{$\mu$-strong convexity}},
  \qquad
  \underbrace{\lVert\nabla f(u)-\nabla f(v)\rVert
  \le L \lVert u-v\rVert}_{\text{$L$-smoothness}},
  \label{eq:strongsmooth}
\end{equation}
with condition number $\kappa = L/\mu$: GD contracts the error by
$(1-1/\kappa)$ per step --- geometry, not luck, sets the speed.
\emph{Convex on a bounded domain} adds existence: a continuous convex $f$
on a compact convex set attains its minimum (Weierstrass), even when
stationarity never fires because the optimum is on the boundary.

\section{Constraints I: equalities and Lagrange multipliers}

Problem: $\min f(\theta)$ s.t.\ $h(\theta)=0$. At a constrained optimum
nothing can be gained by moving \emph{along the constraint surface}, so
$\nabla f$ must be normal to it:
\begin{equation}
  \nabla f(\theta^\ast) + \lambda\, \nabla h(\theta^\ast) = 0,
  \qquad h(\theta^\ast) = 0 ,
  \label{eq:lagrange}
\end{equation}
the stationarity of the \emph{Lagrangian}
$\mathcal{L}(\theta,\lambda) = f(\theta) + \lambda h(\theta)$ in both
arguments. The multiplier is a price: $\lambda = -\,\partial
f^\ast/\partial c$ for the perturbed constraint $h = c$.

\subsection{Walkthrough: the nearest point on a line}

The machinery, executed slowly once (script, panel 2):
\begin{equation*}
  \min_{x,y}\;\; f(x,y) = x^2 + y^2
  \qquad\text{subject to}\qquad
  h(x,y) = x + y - 1 = 0
\end{equation*}
--- geometrically, the point of the line $x+y=1$ closest to the origin.

\paragraph{Step 0: why plain stationarity fails.}
$\nabla f = (2x, 2y)^\top = 0$ only at the origin --- which violates
the constraint. The unconstrained answer is infeasible, so the optimum
sits on the constraint line at a point where $\nabla f \ne 0$, and
Fermat's rule is replaced by \eqref{eq:lagrange}.

\paragraph{Step 1: stationarity in $(x,y)$ --- the shape.}
With $\mathcal{L} = x^2 + y^2 + \lambda(x + y - 1)$:
\begin{equation*}
  \frac{\partial \mathcal{L}}{\partial x} = 2x + \lambda = 0,
  \qquad
  \frac{\partial \mathcal{L}}{\partial y} = 2y + \lambda = 0
  \qquad\Longrightarrow\qquad
  x = y = -\frac{\lambda}{2}.
\end{equation*}
As always, the gradient equations deliver the solution's \emph{shape}
--- here, symmetry: $x = y$, both expressed through the one unknown
multiplier. (The shape is the geometry of \eqref{eq:lagrange}:
$\nabla f = (2x,2y)$ must be parallel to $\nabla h = (1,1)$, which
forces $x = y$.)

\paragraph{Step 2: the constraint fixes the scale.}
One equation remains unused --- the constraint itself. Substitute the
shape into it:
\begin{equation*}
  x + y = -\frac{\lambda}{2} - \frac{\lambda}{2} = -\lambda = 1
  \qquad\Longrightarrow\qquad
  \lambda = -1 .
\end{equation*}

\paragraph{Step 3: back-substitute.}
\begin{equation}
  (x^\ast, y^\ast) = \Bigl(\tfrac12, \tfrac12\Bigr),
  \qquad \lambda = -1 .
  \label{eq:lagexample}
\end{equation}
The perpendicular foot from the origin onto the line --- as the
geometry demanded: at $(\tfrac12,\tfrac12)$ the level circle of $f$ is
tangent to the constraint line, which \emph{is} the parallel-gradients
condition drawn.

\paragraph{The price check.} Perturb the constraint to $x + y = c$: by
the same steps $f^\ast(c) = c^2/2$, so
$-\partial f^\ast/\partial c = -c = -1$ at $c=1$ --- the multiplier's
value, confirming the shadow-price reading. Division of labor to
remember: \textbf{the gradient equations pin the direction of the
answer, the constraint pins its scale, and $\lambda$ is the exchange
rate between objective and constraint} --- solved for, used, discarded.

The multinoulli MLE $\hat p_j = n_j/n$ of the distributions page is this
same machinery applied to $\ell(p) = \sum_j n_j \ln p_j$ under
$\sum_j p_j = 1$ ($\lambda = n$); that instance is walked through at the
same pace in \texttt{lagrange-walkthrough.tex}, next to this note.

\section{Constraints II: inequalities and KKT}

Problem: $\min f(\theta)$ s.t.\ $g_i(\theta) \le 0$, $h_j(\theta) = 0$.

\begin{theorem}[Karush--Kuhn--Tucker conditions]
Under constraint qualification (e.g.\ Slater: some strictly feasible point
exists, for convex problems), a constrained local optimum
$\theta^\ast$ admits multipliers $\mu_i, \lambda_j$ with
\begin{equation}
  \begin{aligned}
    &\text{stationarity:}
    && \nabla f + \textstyle\sum_i \mu_i \nabla g_i
       + \sum_j \lambda_j \nabla h_j = 0 \\
    &\text{primal feasibility:}
    && g_i(\theta^\ast) \le 0, \quad h_j(\theta^\ast) = 0 \\
    &\text{dual feasibility:}
    && \mu_i \ge 0 \\
    &\text{complementary slackness:}
    && \mu_i\, g_i(\theta^\ast) = 0 \quad \text{for all } i.
  \end{aligned}
  \label{eq:kkt}
\end{equation}
For convex $f, g_i$ and affine $h_j$, the KKT conditions are necessary
\emph{and sufficient} for global optimality.
\end{theorem}

Complementary slackness is the operative line: each inequality is either
\emph{inactive} ($g_i < 0$, then $\mu_i = 0$ and it might as well not
exist) or \emph{active} ($g_i = 0$, and $\mu_i \ge 0$ prices it).
Solving a KKT system means guessing the active set, solving the resulting
equality problem, and checking signs --- the script's panel 3 walks the
one-dimensional instance $\min (x-2)^2$ s.t.\ $x \le 1$: the
unconstrained answer $x=2$ is infeasible, the constraint activates,
$x^\ast = 1$, $\mu = 2 > 0$. KKT is the grammar under SVMs (support
vectors are the active constraints), projected methods, and every
`\,s.t.' in an ML paper; Lagrangian \emph{duality} — maximize
$\min_\theta \mathcal{L}$ over the multipliers — is its global companion,
with strong duality under Slater.

\end{document}
